\section{The Diocesan Process and the 1930 Judgment}\label{sec:diocesan}

\subsection{Jurisdiction and process, 1918--1930}

The apparition site lay in the ancient diocese of Leiria, suppressed
in the nineteenth century and restored by Benedict XV in 1918; its
first new bishop, Jos\'e Alves Correia da Silva, took possession in
1920. Competence over the events therefore belonged to him as
diocesan ordinary under the discipline then in force, in which the
judgment of reported apparitions lay with the local bishop---the
regime later codified by the Holy Office's 1978 norms and replaced in
2024 (\S\ref{sec:standing}).\endnote{Diocesan restoration and the bishop's role: the
introductory study and chronology of the \work{Sele\c{c}\~ao de
documentos}; the bishop met Lucia about 1920--1921 (Sanctuary seer
chronology). On the procedural regimes: DDF, \work{Norms for
Proceeding in the Discernment of Alleged Supernatural Phenomena}
(2024), presentation, and its replacement of the norms of 25 February
1978 (art.~27).}

By a \term{provis\~ao} of 3 May 1922 the bishop opened the canonical
process: the document expressly declines to pronounce ``on the origin
and nature of the extraordinary phenomena,'' names a commission of
diocesan clergy to study the alleged facts and organize the process,
orders the text read in every church of the diocese, and counsels the
faithful to ``await the judgment of Holy Church.''\endnote{%
\work{Sele\c{c}\~ao}, Doc.~68 (1922-05-03), publishing DCF~II,
Doc.~1; the commission's character as preparatory is stated in the
commission's own 1930 report (Doc.~120): the provision was issued
``not to pronounce on the origin and nature of the extraordinary
phenomena, but to prepare the elements on which ecclesiastical
authority must base itself in order to render its verdict with full
knowledge of the cause'' (editorial translation).} The process
gathered the 1917 interrogations and testimony; its formal core was
the sworn interrogation of Lucia---by then the only surviving
seer---conducted 8 July 1924 at the Asilo de Vilar in Porto by three
seminary professors with the permission of the Bishop of Porto, the
record being joined to the Leiria process. The commission reported to
the bishop on 13 April 1930.\endnote{\work{Sele\c{c}\~ao}, Doc.~120
(1930-04-13), the commission's report, which recounts the 6 July 1924
request to the Bishop of Porto, the 8 July 1924 sworn interrogation
of the seventeen-year-old Lucia, and the 13 August 1917 detention
episode. Francisco had died 4 April 1919 and Jacinta 20 February 1920
(Sanctuary seer chronology; canonization records, \S\ref{sec:standing}).}

\subsection{The act of 13 October 1930}

The bishop's \work{Carta Pastoral sobre o culto de Nossa Senhora de
F\'atima}, dated at Leiria 13 October 1930 and read at the Cova da
Iria before a great pilgrimage, reviews the events, the hostile and
favorable testimony, the solar phenomenon, the discipline of the
intervening years (including his own long reserve), and the fruits
observed at the site; then, ``having heard the Reverend Consultors of
this our Diocese,'' it issues the two operative clauses quoted at the
head of this study: the visions of the children at the Cova da Iria
on the 13ths of May through October 1917 are declared \emph{worthy of
belief} (\term{dignas de cr\'edito}), and the cult of Our Lady of
F\'atima is officially permitted.\endnote{\work{Sele\c{c}\~ao},
Doc.~133 (1930-10-13), publishing DCF~II, Doc.~11 (also DCF~V-6,
Doc.~1813), printed pages 539--551. The solemn announcement at the Cova
da Iria on the same day, with the bishop's blessing of the multitude,
is documented in the press account printed as Doc.~134 of the same
volume.}

Read strictly, the act does four things and no more. (1)~Its
\emph{object} is ``the visions of the children'' at the Cova da Iria,
May through October 1917: not the 1916 angel apparitions, not the
reported 1915 figures, not Lucia's later experiences at Pontevedra
(1925--1926) or Tuy (1929), not any text of the then-unwritten
memoirs, and not the still-secret parts of the message. (2)~Its
\emph{standard} is credibility---``worthy of belief''---the classical
formula by which an ordinary declares that the faithful may prudently
believe a private revelation, not that they must, and not that its
supernatural origin is a certainty of faith. (3)~It \emph{permits}
public cult, which the diocese had until then only tolerated; this is
the juridical foundation of the Sanctuary. (4)~It binds no
conscience: under the doctrine received then and since, assent to an
approved apparition remains \term{fides humana}, prudent human
credence (\S\ref{sec:obliges}). The act's own argumentative sections---on the solar
phenomenon, on cures, on the movement of pilgrims---are the bishop's
motivation, not further operative determinations.

\subsection{Witness and accessibility ceilings}

Precision about the witness matters. This study quotes the 1930
operative clauses from the Sanctuary of F\'atima's own critical
edition: \work{Documenta\c{c}\~ao Cr\'itica de F\'atima:
Sele\c{c}\~ao de documentos (1917--1930)} (F\'atima: Santu\'ario de
F\'atima, 2013; ISBN 978-972-8213-91-6), Doc.~133, a volume the
Sanctuary distributes without charge as a PDF from its documentation
pages. That edition prints the full Portuguese text with the
editors' apparatus and identifies its earlier publication in the
series (DCF~II, Doc.~11). The ceilings: the original printed pastoral
letter (Uni\~ao Gr\'afica, Lisbon, 1930) and the diocesan archive
original were not consulted; no official English translation of the
act has been identified, so all English here is editorial; and the
fifteen-volume critical series itself was used only through the
Sanctuary's one-volume selection. Within those bounds the wording
rests on the Sanctuary's critical transcription, which is the best
witness generally accessible.\endnote{Edition identity and free
distribution: title, ISBN, and copyright leaf of the
\work{Sele\c{c}\~ao} PDF; download route in the References. The 1930
act is not printed in \work{Acta Apostolicae Sedis}, which records
Holy See acts; a diocesan pastoral letter would not be expected
there, and no AAS locus is claimed.}
